\begin{abstract}
La enseñanza de la teoría musical en el contexto escolar enfrenta importantes desafíos derivados de la reducción de la carga horaria, la discontinuidad curricular y la heterogeneidad en los conocimientos previos de los estudiantes. Aunque existen diversas plataformas digitales de apoyo al aprendizaje musical, estas presentan una limitada capacidad de adaptación pedagógica y carecen de mecanismos que representen explícitamente las relaciones de dependencia entre los contenidos de aprendizaje.

Este trabajo propone un \textbf{modelo de aprendizaje adaptativo} que integra un \textbf{Grafo de Conocimiento Dirigido Acíclico (DAG)} para representar las dependencias pedagógicas entre conceptos de teoría musical y un agente basado en \textbf{Deep Reinforcement Learning (DRL)}, implementado mediante el algoritmo \textit{Proximal Policy Optimization (PPO)}, encargado de recomendar dinámicamente la siguiente actividad de aprendizaje de acuerdo con el estado de conocimiento de cada estudiante. Para facilitar el entrenamiento del agente, el modelo combina una etapa inicial de \textit{Behavior Cloning}, basada en estrategias pedagógicas definidas sobre el DAG, con un proceso posterior de aprendizaje por refuerzo en un entorno de simulación.

El modelo fue implementado en \textbf{SheetMusic}, una plataforma web progresiva orientada al aprendizaje adaptativo de teoría musical y al seguimiento pedagógico de los estudiantes por parte de los docentes. La validación experimental consideró una prueba piloto de 28 días con 28 estudiantes y una evaluación de aceptación tecnológica con docentes de la Región de Los Ríos, Chile.

Los resultados muestran una mejora aproximada del \textbf{20\%} en la mediana del nivel de \textit{proficiency} de los estudiantes, acompañada de una alta percepción de usabilidad (\textbf{SUS = 82.5/100}) y una positiva aceptación tecnológica por parte de los docentes (\textbf{TAM = 3.75/5} en utilidad percibida y \textbf{4.00/5} en facilidad de uso). En conjunto, estos resultados evidencian que la integración de un grafo de conocimiento con \textbf{Deep Reinforcement Learning} constituye una estrategia viable para desarrollar sistemas de aprendizaje adaptativo en el dominio de la teoría musical.


%La educación musical escolar enfrenta severos desafíos derivados de la reducción de la carga horaria, la falta de linealidad curricular y la alta heterogeneidad en los conocimientos de los estudiantes. Las herramientas digitales existentes carecen de alineación curricular local, flexibilidad adaptativa y mecanismos de gestión pedagógica para el cuerpo docente. Este trabajo presenta el diseño, implementación y validación de SheetMusic, una plataforma web progresiva y servicio web distribuido para la enseñanza adaptativa de la teoría musical. El sistema articula la experiencia de aprendizaje del estudiante mediante un grafo de conocimiento no lineal y gamificado, herramientas de seguimiento pedagógico para el docente, y un microservicio inteligente basado en Aprendizaje por Refuerzo Profundo (DRL) mediante el algoritmo Proximal Policy Optimization (PPO) que ajusta en tiempo real la dificultad de las actividades. El conocimiento musical se formalizó en un Grafo Dirigido Acíclico (DAG) de 23 nodos que modela las dependencias pedagógicas. La validación experimental abarcó pruebas piloto durante 28 días con 28 estudiantes y un estudio cuali-cuantitativo (entrevistas y pruebas de usabilidad) con 6 docentes de la Región de Los Ríos, Chile. Los resultados demostraron una mejora en la mediana de desempeño (proficiency) del 20\% impulsada por el motor DRL, una alta usabilidad percibida por parte de los estudiantes (SUS de 82.5/100) y una positiva aceptación tecnológica docente (TAM con 3.75/5 en utilidad percibida y 4.00/5 en facilidad de uso).
\end{abstract}

\begin{IEEEkeywords}
Aprendizaje Adaptativo, Deep Reinforcement Learning, Grafo de Conocimiento, Recomendación Adaptativa, Teoría Musical, Sistemas Inteligentes Educativos.
%Aprendizaje Adaptativo, Educación Musical, Servicio Web, Deep Reinforcement Learning, Grafos de Conocimiento, Usabilidad, TAM.
\end{IEEEkeywords}